\section{Related Work}
\label{sec:related}

\para{Formalizing vibes in NLP.}
\citet{dunlap2024vibecheck} introduced \textsc{VibeCheck}, a system that discovers and quantifies qualitative ``vibe'' axes---tone, formatting, writing style---along which LLM outputs differ.
Their pipeline uses an LLM to examine output pairs, discover vibe axes (e.g., ``formal $\to$ friendly''), and score outputs on each axis.
While \textsc{VibeCheck} demonstrates that vibes are measurable and predictive (achieving 80\% accuracy at identifying model identity), it is designed for comparing LLM outputs, not for characterizing web content or constructing navigable embedding spaces.
We extend the concept of vibes from model comparison to web-page description and embed the resulting descriptions to create a spatial map.

\para{Embedding visualization.}
Interactive visualization of high-dimensional embeddings is a well-studied problem.
\citet{smilkov2016embedding} introduced the Embedding Projector for exploring word and image embeddings via t-SNE, UMAP, and PCA.
\citet{polo2023wizmap} proposed \textsc{WizMap}, which scales to millions of points using multi-resolution map-like navigation.
Most recently, \citet{ren2025embedding} presented Embedding Atlas, a browser-based tool that supports real-time UMAP computation, density contours, and cross-filtering on millions of data points.
All of these tools visualize \emph{content} or \emph{semantic} embeddings; none address the question of whether a different embedding---one capturing aesthetic character rather than topic---would produce a qualitatively different and useful map.
Our work contributes a new kind of input to these visualization tools: vibe embeddings.

\para{Interpreting embedding spaces.}
\citet{tennenholtz2023demystifying} proposed ELM (Embedding Language Model), which trains adapter layers that allow an LLM to generate natural-language descriptions of entities from their embeddings.
\citet{liu2024visualizing} developed gradient-based methods to overlay spatial word clouds on 2D document projections, making regions of embedding space interpretable.
\citet{haschka2025multiscale} build hierarchical semantic trees from embedding spaces, revealing multi-scale cluster structure.
These approaches interpret \emph{existing} embedding spaces; we construct a \emph{new} embedding space by first translating content into vibe descriptions and then embedding those descriptions.

\para{Style and aesthetic representation.}
Recent work has examined how writing style, as distinct from topic, shapes the geometry of embedding spaces~\cite{embedding_style2025}.
\citet{navigation2026} model human concept navigation as trajectories through transformer embedding spaces, providing a theoretical framework for browsing semantic maps.
Our work is complementary: we show that vibe descriptions---which capture style, era, community type, and cultural energy---produce embeddings that organize content along aesthetic axes rather than topical ones.

\para{LLMs and informal language.}
The ``it's giving'' prompt we use draws on contemporary internet slang.
\citet{slang2024} and \citet{wilson2024slang} study how LLMs comprehend and produce slang, finding that models trained on large web corpora handle informal language reasonably well.
\citet{gur2022understanding} demonstrate that LLMs pretrained on natural language transfer effectively to HTML understanding tasks.
We build on these findings by prompting an LLM with web page text and a slang-inflected question, relying on the model's cultural fluency to produce evocative, aesthetic descriptions.
